\documentclass[UTF8,11pt,a4paper,fontset=fandol]{ctexart}
\usepackage[left=24mm,right=24mm,top=23mm,bottom=23mm,headheight=15pt]{geometry}
\usepackage{amsmath,amssymb,mathtools,bm}
\usepackage{xcolor,fancyhdr,needspace}
\usepackage[unicode,colorlinks=true,linkcolor=black,urlcolor=blue!45!black]{hyperref}
\definecolor{ink}{HTML}{183B4E}
\hypersetup{pdftitle={从库仑交换到二维铁磁体的单磁子色散},pdfauthor={},pdfsubject={自旋二分之一铁磁体的单翻转运动方程与色散}}
\pagestyle{fancy}
\fancyhf{}
\fancyhead[L]{\small\color{ink}二维铁磁体的单磁子色散}
\fancyfoot[C]{\small\thepage}
\renewcommand{\headrulewidth}{0.3pt}
\setlength{\parindent}{2em}
\setlength{\parskip}{0.3em}
\linespread{1.18}
\allowdisplaybreaks[1]
\newcommand{\ket}[1]{\lvert #1\rangle}
\newcommand{\bra}[1]{\langle #1\rvert}
\newcommand{\id}{\mathbb I}
\begin{document}
\begin{center}
{\LARGE\bfseries\color{ink}从库仑交换到二维铁磁体的单磁子色散}\\[0.65em]
{\small 两电子矩阵、单翻转子空间与运动方程}
\end{center}
\vspace{0.3em}

两电子之间的库仑相互作用为
\begin{equation}
v(\mathbf r_1-\mathbf r_2)=
\frac{e^2}{4\pi\varepsilon_0|\mathbf r_1-\mathbf r_2|},
\end{equation}
其二次量子化形式为
\begin{equation}
\hat V=\frac12\sum_{\sigma,\tau}
\int d^3r_1\,d^3r_2\,
\hat\psi^\dagger_\sigma(\mathbf r_1)
\hat\psi^\dagger_\tau(\mathbf r_2)
v(\mathbf r_1-\mathbf r_2)
\hat\psi_\tau(\mathbf r_2)\hat\psi_\sigma(\mathbf r_1).
\end{equation}
取两个正交归一轨道 $a,b$，各占据一个电子，场算符展开为
$\hat\psi_\sigma(\mathbf r)=\phi_a(\mathbf r)f_{a\sigma}+\phi_b(\mathbf r)f_{b\sigma}$。
在基底
\begin{equation}
\bigl\{
\ket{\uparrow_a\uparrow_b},\
\ket{\uparrow_a\downarrow_b},\
\ket{\downarrow_a\uparrow_b},\
\ket{\downarrow_a\downarrow_b}
\bigr\},\qquad
\ket{\sigma_a\tau_b}=f^\dagger_{a\sigma}f^\dagger_{b\tau}\ket{0}
\label{eq:pairbasis}
\end{equation}
下，$\hat V$ 的矩阵为
\begin{equation}
V_{ab}=
\begin{pmatrix}
D-K&0&0&0\\
0&D&-K&0\\
0&-K&D&0\\
0&0&0&D-K
\end{pmatrix},
\label{eq:coulombmatrix}
\end{equation}
其中直接积分与交换积分分别为
\begin{align}
D&=\int d^3r_1\,d^3r_2\,
|\phi_a(\mathbf r_1)|^2\,v(\mathbf r_1-\mathbf r_2)\,
|\phi_b(\mathbf r_2)|^2, \label{eq:direct}\\
K&=\int d^3r_1\,d^3r_2\,
\phi_a^*(\mathbf r_1)\phi_b^*(\mathbf r_2)\,
v(\mathbf r_1-\mathbf r_2)\,
\phi_b(\mathbf r_1)\phi_a(\mathbf r_2). \label{eq:exchange}
\end{align}

取 $\hbar=1$，以 $x$ 轴为自旋量子化轴，$\uparrow,\downarrow$ 分别表示沿 $+x,-x$ 的自旋态。在同一基底下，
\begin{equation}
\mathbf S_a\cdot\mathbf S_b
=\frac14
\begin{pmatrix}
1&0&0&0\\
0&-1&2&0\\
0&2&-1&0\\
0&0&0&1
\end{pmatrix},
\label{eq:dotmatrix}
\end{equation}
因而
\begin{equation}
\boxed{
V_{ab}=\left(D-\frac K2\right)\id_4
-2K\,\mathbf S_a\cdot\mathbf S_b.
}
\label{eq:pairspin}
\end{equation}

对最近邻键求和，略去常数项，令 $J_{ij}=2K_{ij}$，得到海森堡哈密顿量
\begin{equation}
H=-\sum_{\langle ij\rangle}J_{ij}\,\mathbf S_i\cdot\mathbf S_j.
\label{eq:latticeH}
\end{equation}
对于无外场的二维正方晶格，取水平、竖直方向的铁磁交换分别为 $J_x,J_y>0$，则
\begin{equation}
H=-J_x\sum_{\langle ij\rangle_x}\mathbf S_i\cdot\mathbf S_j
-J_y\sum_{\langle ij\rangle_y}\mathbf S_i\cdot\mathbf S_j.
\label{eq:squareH}
\end{equation}
每条键计一次，晶格常数为 $a$，格点数为 $\mathcal N$，采用周期边界。

全极化基态为
\begin{equation}
\ket{F}=\ket{\uparrow_1\uparrow_2\cdots\uparrow_{\mathcal N}}.
\end{equation}
每条键上 $\mathbf S_i\cdot\mathbf S_j$ 的值为 $1/4$，两个方向各有 $\mathcal N$ 条键，因此
\begin{equation}
H\ket{F}=E_F\ket{F},\qquad
\boxed{E_F=-\frac{\mathcal N}{4}(J_x+J_y).}
\label{eq:EF}
\end{equation}

单翻转子空间的基底为
\begin{equation}
\ket{i}=
\ket{\uparrow_1\cdots\uparrow_{i-1}\downarrow_i
\uparrow_{i+1}\cdots\uparrow_{\mathcal N}}.
\label{eq:oneflip}
\end{equation}
对连接 $i$ 与其邻居 $j$ 的键，将 $\ket{\downarrow_i\uparrow_j}$ 代入矩阵 \eqref{eq:dotmatrix}，得到
\begin{equation}
-J_{ij}\mathbf S_i\cdot\mathbf S_j\ket{i}
=\frac{J_{ij}}4\ket{i}-\frac{J_{ij}}2\ket{j}.
\label{eq:bondaction}
\end{equation}
其余键的贡献与全极化态相同。对所有键求和，
\begin{equation}
H\ket{i}=E_F\ket{i}
+\frac12\sum_{j\text{ 邻接 }i}J_{ij}
\bigl(\ket{i}-\ket{j}\bigr).
\label{eq:Haction}
\end{equation}
因此单翻转子空间在时间演化下封闭。初始只翻转格点 $i_0$，写成
\begin{equation}
\ket{\Psi(t)}=\sum_i c_i(t)\ket{i},\qquad
c_i(0)=\delta_{i,i_0}.
\end{equation}
代入薛定谔方程 $i\partial_t\ket{\Psi}=H\ket{\Psi}$，比较各 $\ket{i}$ 的系数，得到
\begin{equation}
\boxed{
i\dot c_i
=E_Fc_i+\frac12\sum_{j\text{ 邻接 }i}J_{ij}(c_i-c_j).
}
\label{eq:EOMgeneral}
\end{equation}
写出正方晶格的四个邻居，
\begin{equation}
\boxed{
\begin{aligned}
i\dot c_i={}&(E_F+J_x+J_y)c_i\\
&-\frac{J_x}{2}
(c_{i+\mathbf e_x}+c_{i-\mathbf e_x})
-\frac{J_y}{2}
(c_{i+\mathbf e_y}+c_{i-\mathbf e_y}),
\end{aligned}
}
\label{eq:EOMsquare}
\end{equation}
其中 $i\pm\mathbf e_x,i\pm\mathbf e_y$ 表示水平、竖直方向的相邻格点。

取平面波模式
\begin{equation}
c_i(t)=\frac1{\sqrt{\mathcal N}}
\exp\!\left[i\mathbf k\cdot\mathbf r_i-iE_{\mathbf k}t\right],
\label{eq:planewave}
\end{equation}
其中 $\mathbf r_i$ 为格点位置，$\mathbf k=(k_x,k_y)$ 为周期边界允许的波矢。将
\begin{equation}
c_{i\pm\mathbf e_x}=e^{\pm ik_xa}c_i,\qquad
c_{i\pm\mathbf e_y}=e^{\pm ik_ya}c_i,\qquad
i\dot c_i=E_{\mathbf k}c_i
\end{equation}
代入 \eqref{eq:EOMsquare}，得到
\begin{equation}
\begin{aligned}
E_{\mathbf k}
&=E_F+J_x+J_y
-\frac{J_x}{2}(e^{ik_xa}+e^{-ik_xa})
-\frac{J_y}{2}(e^{ik_ya}+e^{-ik_ya})\\
&=E_F+J_x(1-\cos k_xa)+J_y(1-\cos k_ya).
\end{aligned}
\end{equation}
相对全极化基态的单磁子激发能为
\begin{equation}
\boxed{
E_{\mathbf k}-E_F
=J_x(1-\cos k_xa)+J_y(1-\cos k_ya).
}
\end{equation}
\end{document}
